\section{Formación básica de backend}

\subsection{Qué es el backend}
El backend es la parte que procesa la información y las reglas de AUTOVIDA fuera de la pantalla que ve la persona usuaria. De forma simplificada, la interfaz envía una petición; el servidor comprueba que sea válida y que la persona tenga permiso; consulta o actualiza los datos almacenados; y devuelve un resultado que la interfaz presenta de forma comprensible.

Por ejemplo, al guardar una tarea, el servidor debe comprobar que los campos requeridos tienen valores válidos, que quien la guarda puede modificarla, y que la operación se ha almacenado correctamente. La interfaz puede ayudar a prevenir errores, pero nunca debe ser la única barrera: los datos recibidos por el servidor se consideran no confiables hasta validarlos.

\subsection{Conceptos esenciales}
\begin{itemize}[leftmargin=*]
  \item \textbf{API:} conjunto definido de operaciones por las que la interfaz y el servidor intercambian datos. Cada operación debe tener entradas y respuestas documentadas.
  \item \textbf{Base de datos:} almacenamiento organizado de la información. Debe evitar registros inconsistentes y guardar solo los datos necesarios para las funciones acordadas.
  \item \textbf{Autenticación y autorización:} la autenticación identifica quién accede; la autorización determina qué acciones puede realizar esa persona. Son comprobaciones distintas.
  \item \textbf{Validación:} reglas aplicadas a los datos recibidos, como campos obligatorios, formatos y límites permitidos. Deben comprobarse en el servidor.
  \item \textbf{Respuesta y error:} resultado que el servidor comunica a la interfaz. Debe ser consistente y permitir mostrar una explicación útil sin revelar información interna.
\end{itemize}

\subsection{Flujo que debe entender todo el equipo}
Para cada función, dibujar o describir estos pasos: acción del usuario, petición enviada, validaciones, regla de negocio, lectura o escritura de datos, respuesta y mensaje visible. Este esquema ayuda a detectar requisitos olvidados y a planificar pruebas aunque no todas las personas programen backend.

\begin{infobox}[title={Idea clave}]
El frontend facilita una interacción comprensible; el backend protege las reglas y los datos. Nunca se debe confiar en que la interfaz haya enviado información correcta o de que una persona tenga permiso solo porque un botón estaba oculto.
\end{infobox}